\section{Milk spectra construction and analysis}\label{supp:milk}
\newcommand{\MilkCows}{1349}
\newcommand{\MilkSamples}{2253}
\newcommand{\MilkRepeatCows}{904}
\newcommand{\MilkRepeatSamples}{1808}
\newcommand{\MilkOriginalCows}{1356}
\newcommand{\MilkOriginalSamples}{2270}
\newcommand{\MilkInvalidCows}{7}
\newcommand{\MilkInvalidSamples}{17}
\newcommand{\MilkSpectra}{1060}
\newcommand{\MilkDimension}{1073}
\newcommand{\MilkAOneCows}{519}
\newcommand{\MilkATwoCows}{830}
\newcommand{\MilkFits}{18}
\newcommand{\MilkSystems}{72}
\newcommand{\MilkMinSelected}{27}
\newcommand{\MilkMaxSelected}{163}
\newcommand{\MilkMedianRuntime}{17.0}
\newcommand{\MilkMaxKKT}{5.17e-09}
\newcommand{\MilkMaxNuisance}{9.99e-11}
\newcommand{\MilkSpectralCorrelation}{0.9986}
\newcommand{\MilkLowerEstimate}{-0.020}
\newcommand{\MilkLowerSE}{0.101}
\newcommand{\MilkLowerInterval}{[-0.263,\ 0.223]}
\newcommand{\MilkMedianEstimate}{0.018}
\newcommand{\MilkMedianSE}{0.118}
\newcommand{\MilkMedianInterval}{[-0.264,\ 0.299]}
\newcommand{\MilkUpperEstimate}{0.119}
\newcommand{\MilkUpperSE}{0.140}
\newcommand{\MilkUpperInterval}{[-0.216,\ 0.453]}


\subsection{Recorded measurements and cow identifiers}

The dataset \cite{chirife2025data} contains \MilkOriginalSamples{}
milk samples from \MilkOriginalCows{} cow identifiers. The Excel
columns numbered 240 through 1299 contain the \MilkSpectra{} MIR
absorbance measurements. Somatic cell count is recorded in cells
per microlitre and transformed as $\log_2(\mathrm{SCC}/100)+3$.
All retained spectral values and counts are finite, and every count
is positive.

Eligibility requires each cow identifier to have a consistent farm
and milk-protein genotype, with one record per recorded milk control
and distinct sampling dates for repeated observations. This rule
excludes \MilkInvalidSamples{} records belonging to
\MilkInvalidCows{} identifiers and retains \MilkSamples{} samples
from \MilkCows{} cows. The original records are preserved. The
paired cohort contains \MilkRepeatSamples{} observations from
\MilkRepeatCows{} cows. Recorded sampling dates span 24 May through
31 July 2023, and the intervals between repeated measurements range
from 28 to 66 days.

The milk-type coordinate is the recorded A2 indicator, equal to one
for A2A2 and zero for A1A1 or A1A2. Its definition agrees with the
deposited milk-type field for every retained observation. All six
farms contain both indicator values. The twelve farm and milk-control
indicators and the milk-type indicator have full column rank in
each outer selection sample and its inner training samples.

Table~\ref{tab:milk-descriptive} summarises the observed score
distributions. Means and empirical quantiles give each cow total
weight one by assigning weight $1/m_i$ to each of its records.

\begin{table}[t]
\centering
\small
\caption{Observed somatic cell score distributions with equal total weight per cow. Each record receives weight $1/m_i$. The type-specific summaries average over the recorded farm conditions and spectra.}
\label{tab:milk-descriptive}
\begin{tabular}{lrrrrrr}
\toprule
Cows & Number & Samples & Mean & Lower quartile & Median & Upper quartile \\
\midrule
All eligible cows & 1349 & 2253 & 2.723 & 1.214 & 2.456 & 3.880 \\
A1 milk type & 519 & 871 & 2.776 & 1.310 & 2.506 & 3.949 \\
A2 milk type & 830 & 1382 & 2.690 & 1.214 & 2.401 & 3.850 \\
One milk control & 445 & 445 & 2.901 & 1.401 & 2.546 & 4.021 \\
Two milk controls & 904 & 1808 & 2.636 & 1.163 & 2.356 & 3.799 \\
\bottomrule
\end{tabular}
\end{table}


\subsection{Sample roles and the high dimensional design}

Seed 20261031 generates a random ordering of the retained cow
identifiers. Alternating positions in that ordering define the
two complementary groups. Each group performs selection for the
other group's coefficient estimation. Inner selection and validation
roles also contain complete cows. The paired cohort applies the
same deterministic fold rule to its retained identifiers.

The full design has \MilkDimension{} columns. The spectral block
enters the outcome and adjustment searches at its recorded resolution.
RMS scales in the Gaussian and outcome-search subproblems are
calculated from their own training observations. The final milk-type
coefficient is expressed for a change from indicator zero to one.
All 13 prespecified coordinates remain included in the selected
profile. Outcome search size and graph penalties follow the method's
cluster validation rules; their losses and selected supports are
saved for each estimation group.

The outcome search and curvature graphs depend on the quantile and
working parameters. The conditional design regression uses only the
covariates and cluster identifiers. Its four combinations of cohort
and estimation group are computed once and shared across settings.
Each reuse verifies the complete design matrix and the ordered cow
identifiers. Gaussian Lasso paths use the covariance algorithm in
\texttt{glmnet}, with convergence threshold $10^{-13}$ and a maximum
of $10^7$ iterations. Every path point must have KKT residual at most
$10^{-6}$. The ratio grids and validation losses retain the
definitions in the method specification.

\subsection{Equivalent coordinates for the selected fit}

Let $W_0$ be diagonal with entries $1/(n_fm_i)$ on an estimation
sample, and write the weighted selected design as
$W_0^{1/2}X_AP=QR$, where $P$ is the QR column permutation.
The design $U=W_0^{-1/2}Q$ and coefficient
$a=RP^T\beta_A$ satisfy $X_A\beta_A=Ua$. Moreover,
\[
 \beta_A^T(X_A^TW_0X_A)\beta_A=a^Ta.
\]
Thus the prediction ridge and complete profile objective are
preserved. Newton iterations use $a$, and the reported coefficient
is $\beta_A=PR^{-1}a$.

For a target coordinate $k$, define
$c_k=R^{-T}P^Te_k$. If $H_a$ is the fitted curvature including the
prediction ridge in the transformed coordinates, the original
inverse direction is $PR^{-1}H_a^{-1}c_k$. Cluster projected scores
are obtained by multiplying $-U_iH_a^{-1}c_k$ elementwise by the
within-cow scores and averaging over that cow. This computes the
same inverse row and influence quantity as the original selected
criterion. The numerical checks compare original-coordinate gradients
and linear predictions. An actual-data reference calculation also
compares coefficients, objectives, Hessians and projected cluster
scores under the two coordinate representations.

\subsection{Intervals and sensitivity calculations}

For each quantile, the estimate combines both groups' paired
coefficients with their cow-count weights. The standard error is
$n^{-1}(\sum_f\sum_{i\in I_f}(V_{ik}-\bar V_{fk})^2)^{1/2}$.
The saved influence values reproduce this quantity to numerical
precision. Pointwise intervals use $z_{0.975}$; simultaneous
intervals use $z_{1-0.05/6}$ for the three quantiles within a setting.

\subsection{Supplementary comparisons across quantiles}

Write $\widehat\theta=(\widehat\theta_{0.25},
\widehat\theta_{0.50},\widehat\theta_{0.75})^T$ and collect the
corresponding saved cow influence contributions into the vector
$V_i$. Cow identifiers align these contributions across quantiles.
The joint covariance estimate is
\[
 \widehat\Sigma=\frac1{n^2}\sum_f\sum_{i\in I_f}
 (V_i-\bar V_f)(V_i-\bar V_f)^T.
\]
Its diagonal reproduces the three reported coordinate variances.
For a contrast $d$, the estimate is $d^T\widehat\theta$ and its
standard error is $(d^T\widehat\Sigma d)^{1/2}$. The three rows
$(-1,1,0)$, $(-1,0,1)$ and $(0,-1,1)$ define all pairwise differences.
Their intervals use a separate Bonferroni family of three.
This supplementary calculation uses the saved reference fits.

For the joint equality hypothesis, let $G$ contain the first two
contrast rows. The Wald statistic is
\[
 W=(G\widehat\theta)^T(G\widehat\Sigma G^T)^{-1}G\widehat\theta.
\]
The finite vector of influence expansions gives a joint normal
approximation under the stated moment and learning conditions.
With nonsingular limiting contrast covariance, the null reference
distribution is $\chi^2_2$. The observed statistic is
$\MilkQuantileEqualityWald{}$, giving $p=\MilkQuantileEqualityP{}$.
The upper-minus-lower contrast is $\MilkUpperLowerDifference{}$,
with simultaneous interval $\MilkUpperLowerDifferenceInterval{}$.

\begin{table}[t]
\centering
\small
\caption{Supplementary differences between the three reference profile coefficients. Standard errors retain their common cow contributions. The simultaneous intervals and adjusted $p$ values use a Bonferroni family of three differences.}
\label{tab:milk-contrasts}
\begin{tabular}{lrrrr}
\toprule
Difference & Estimate & Standard error & Simultaneous interval & Adjusted $p$ \\
\midrule
$\theta_{0.50}-\theta_{0.25}$ & $0.038$ & 0.116 & $[-0.241,\ 0.317]$ & 1.000 \\
$\theta_{0.75}-\theta_{0.25}$ & $0.139$ & 0.143 & $[-0.204,\ 0.482]$ & 0.996 \\
$\theta_{0.75}-\theta_{0.50}$ & $0.101$ & 0.143 & $[-0.241,\ 0.443]$ & 1.000 \\
\bottomrule
\end{tabular}
\end{table}


\subsection{Sensitivity to the analysis settings}

Table~\ref{tab:milk-sensitivity} includes every registered setting.
The two-control cohort has its own sampling population, and each
working penalty is part of its profile parameter definition.
Bandwidth sensitivity retains the same population target and
changes the numerical smoothing scale.

\begin{table}[p]
\centering
\small
\caption{Prespecified sensitivity analyses. Each setting reports its own population profile parameter. Intervals allocate 5\% across its three quantiles. Selected dimensions are the minimum and maximum over the two estimation samples. Time is the complete coordinate-fitting time in seconds after the shared design regression has been computed.}
\label{tab:milk-sensitivity}
\begin{tabular}{lrrrrr}
\toprule
Setting & $\tau$ & Estimate & Simultaneous interval & Dimension & Seconds \\
\midrule
Reference & 0.25 & -0.020 & $[-0.263,\ 0.223]$ & 29--90 & 21.2 \\
Reference & 0.50 & 0.018 & $[-0.264,\ 0.299]$ & 100--124 & 26.7 \\
Reference & 0.75 & 0.119 & $[-0.216,\ 0.453]$ & 32--146 & 16.5 \\
Two milk controls & 0.25 & -0.137 & $[-0.428,\ 0.155]$ & 28--32 & 10.3 \\
Two milk controls & 0.50 & -0.248 & $[-0.590,\ 0.094]$ & 29--43 & 11.6 \\
Two milk controls & 0.75 & -0.205 & $[-0.697,\ 0.288]$ & 45--91 & 9.9 \\
$0.75h$ & 0.25 & -0.019 & $[-0.276,\ 0.237]$ & 29--90 & 21.7 \\
$0.75h$ & 0.50 & 0.049 & $[-0.250,\ 0.348]$ & 105--124 & 23.4 \\
$0.75h$ & 0.75 & 0.114 & $[-0.236,\ 0.465]$ & 33--148 & 11.9 \\
$1.25h$ & 0.25 & -0.021 & $[-0.274,\ 0.231]$ & 29--90 & 15.0 \\
$1.25h$ & 0.50 & 0.016 & $[-0.278,\ 0.310]$ & 100--124 & 16.9 \\
$1.25h$ & 0.75 & 0.107 & $[-0.245,\ 0.460]$ & 34--148 & 12.7 \\
$\Lambda=0.5$ & 0.25 & -0.023 & $[-0.249,\ 0.203]$ & 27--87 & 18.6 \\
$\Lambda=0.5$ & 0.50 & 0.072 & $[-0.199,\ 0.344]$ & 85--141 & 21.0 \\
$\Lambda=0.5$ & 0.75 & 0.176 & $[-0.134,\ 0.485]$ & 142--163 & 19.0 \\
$\Lambda=2$ & 0.25 & -0.160 & $[-0.373,\ 0.053]$ & 38--89 & 16.0 \\
$\Lambda=2$ & 0.50 & 0.005 & $[-0.332,\ 0.341]$ & 31--131 & 17.2 \\
$\Lambda=2$ & 0.75 & 0.173 & $[-0.371,\ 0.716]$ & 81--150 & 18.6 \\
\bottomrule
\end{tabular}
\end{table}


Across the \MilkSystems{} coefficient systems, the maximum
original-coordinate KKT residual is \texttt{\MilkMaxKKT{}} and the
maximum nuisance gradient is \texttt{\MilkMaxNuisance{}}.
Selection roles, every chosen variable, curvature diagnostics and
all fitted cluster influence values are retained with the
application output. Each completed case is copied to the separate
backup directory before the queue advances.

\bibliography{references}
\end{document}
